\chapter*{Preface}
\textbf{I have zero formal education in electronics}. I build
circuits as a hobby and have fallen in love with the subject.

I wrote this notebook to document the mathematical analysis
of various electrical circuits I encounter in the wild. My aim
is to keep theory to a minimum and therefore to utilise this
notebook effectively, you must be familiar with circuit theory,
control systems, and linear algebra.

If you seek a more formal and comprehensive understanding of
electronics, I recommend the following textbooks:
\begin{enumerate}
    \item \textit{The Art of Electronics} by Paul Horowitz and Winfield Hill
    \item \textit{Practical Electronics for Inventors} by Paul Scherz and Simon Monk
\end{enumerate}

I want to mention my gratitude to the authors of the above
textbooks, as they have been instrumental in my understanding
of electronics. I also want to thank countless online resources
and professors who have shared their knowledge and expertise in
the field.

This \textbf{notebook is a work in progress}, I welcome feedback
and suggestions for improvement. If you find any errors or have
ideas for new topics to cover, please raise an issue on the GitHub
repository or reach out to me directly.
\\
Best regards,\\
Ayush Gupta

\chapter*{Organisation}
\addcontentsline{toc}{chapter}{Organisation}
Every block is a combination of devices. Every block is analysed with
the octant system of Chapter~\ref{ch:analysis}. The notebook therefore
reads bottom up:
\begin{verbatim}
 Analysis ......... octant system (how to analyse anything below)
 Devices .......... Passives: Resistor, Capacitor, Inductor, Diode
                    Actives:  BJT, MOSFET, Op-Amp
 General .......... Voltage Dividers, Current Dividers
 Filters .......... Passive: Single Pole (RC, RL), Double Pole (LC, RLC)
                    Active:  Single Pole, Sallen-Key, Multiple Feedback
 Current Mirrors .. Basic, Widlar, Wilson, Modified Wilson, Cascode
 Amplifiers ....... Common:       CS, CG, CD, CE, CB, CC
                    Functional:   Buffer, Inverting, Non-Inverting,
                                  Summing, Difference, Integrator,
                                  Differentiator
                    Differential: Differential Pair, 5T OTA, Cascode,
                                  Telescopic Cascode, Regulated Cascode,
                                  Folded Cascode, Folded Regulated
                                  Cascode, Current Mirror Op-Amp,
                                  Two-Stage Op-Amp
 Data Converters .. Characterisation
                    DAC: String, R-2R, Current Source/Sink, MDAC,
                         Current Steering (1-Steering), SAR, Pipeline,
                         Delta-Sigma
                    ADC: Nyquist Rate: Flash, Two-Step Flash, Multi-Step
                         Flash, Pipeline, Cyclic, Interpolating,
                         Folding, SAR, Dual Slope
                         Over-Sampled: First Order, Higher Order,
                         Single-Bit, Multi-Bit, Continuous-Time
 PLLs and VCOs .... Phase Detectors, Charge Pump, Loop Filter, PLL,
                    Ring VCO, LC VCO, Relaxation VCO
 RF ............... (reserved)
 Electromechanics . Motors
 Appendix ......... Theory: General, Controls, Electromagnetics
                    Symbols and Notation
\end{verbatim}
Each block follows the same layout: \textbf{Overview}, \textbf{Schematic},
\textbf{Transfer Function} (Large Signal, Time Domain, DC Input Analysis
(Q5), and Small Signal, Frequency Domain, AC Input Analysis (Q4)),
\textbf{Analysis} (poles, zeros, gains, impedances), \textbf{Design},
\textbf{Question} and, where the algebra is heavy, \textbf{Code}.

Notation and every symbol and label used in the notebook are collected
in Appendix~\ref{ch:symbols}.
